%% Anotações da lição 04 — exemplo com hebraico e numeração dos versículos
\EBlivcampo{Ideia central}{a resposta da lição, em uma ou duas frases}

\begin{campo}{Texto em foco: Salmo 83:18}
\begin{alerta}[Numeração]
Nas Bíblias em português o texto é \rb{Sl 83:18}; no hebraico é
\rb{83:19}, porque o título do salmo conta como versículo~1. Ao gerar dados
com \texttt{interlinear.py}, use a numeração hebraica.
\end{alerta}

\interlinear[titulo={Sl 83:19 (hebraico) · WLC}]{dados/biblia/sl083_19.hbo.tsv}

Observe a construção: \hb{כִּי אַתָּה שִׁמְךָ יְהוָה לְבַדֶּךָ עֶלְיוֹן}. As
traduções diferem em como ligam \enquote{teu nome} e \enquote{só tu}
(\enquote{que tu, cujo nome é\ldots, só tu és o Altíssimo} ou \enquote{que só
tu, cujo nome é\ldots, és o Altíssimo}). Compare duas traduções e anote.
\end{campo}

\EBlivcampo[3]{Dúvidas para pesquisar}{}
\EBlivcampo[3]{Revisão sem consultar}{responda às perguntas de revisão da lição de memória}
